\chapter{Introduction}
The economical and safety related reasons had spurred ever-increasing interest in the theory and applications of
model-based fault detection and diagnosis methods. The use of parameter estimation techniques for fault detection of
mechanical systems was first demonstrated in \cite{baskiotis1979parameter}, \cite{geiger1982monitoring}
\cite{filbert1982quality}, \cite{hohmann1977automatic}. This as a part of larger class of fault detection based on
analytical redundancy. The development of process fault detection methods based on modelling, parameter and state
estimation was summarized by \cite{isermann1984process}, \cite{isermann1997trends}. Parity equation-based based methods
were treated early in \cite{chow1984analytical} and then later developed by Patton and others in \cite{patton1991review}
and \cite{hofling1994detection}.

Basic terminology related to fault diagnosis can be found in \cite{hofling1994detection} and \cite{ding2005model}. The following terms need the specific attention:

\begin{itemize}
        \item \itbf{Fault Detection}: Determination of faults present in the system at the time of detection.
        \item \itbf{Fault Isolation}: Determination of the kind, location and time of detection of a fault. Follows fault detection.
        \item \itbf{Fault Identification}: Determination of the size and time variant behavior of the fault. Follows fault isolation.
        \item \itbf{Fault Diagnosis}: Determination of the kind, size, location and time of detection of a fault. Follows fault detection. Includes fault detection and identification.
        \item \itbf{Residual}: A fault indicator signal, based on the deviation between the measurements and model equation based computations.
        \item \itbf{Analytical redundancy}: Use of more (not necessarily identical) ways to determine a variable, where
                one way uses a mathematical process model in analytical form. Analytical or functional redundancy
                exploits redundant analytical relationships among various measured variables of the monitored process \cite{chen2012robust}.
        \item \itbf{Fault Classification terminology}:
        \begin{itemize}
                \item \itbf{Time dependency of fault}:
                \begin{itemize}
                        \item \itbf{Abrupt Fault}: Fault modelled as stepwise function. It represents bias in the monitored signal.
                        \item \itbf{Incipient fault}: Fault modelled by using ramp signals. It represents drift of the monitored signal.
                        \item \itbf{Intermittent fault}: Combination of impulses with different amplitudes.
                \end{itemize}
                \item \itbf{Analytical type of fault}:
                \begin{itemize}
                        \item \itbf{Analytical redundancy}: Influences a variable by an addition of the fault itself. They may represent, e.g., offsets of sensors.
                        \item \itbf{Multiplicative fault}: Are represented by the product of a variable with the fault itself. They can appear as parameter changes within a process.
                \end{itemize}
        \end{itemize}
\end{itemize}


Consistency checking in analytical redundancy is normally achieved through a comparison between a measured signal with
estimated values. The estimation is generated by a mathematical model of the considered plant. The comparison is done
using the residual quantities which are computed as differences between the measured signals and the corresponding
signals generated by the mathematical model. Observers, Parity equations and Identification (Parameter Estimation) are
some basic model based residual generation methods. The characteristic quantities or features from fault detection
methods show stochastic behavior with mean values and variances. Deviations from the normal behavior must then be
detected by methods of change detection.

%=======================================================================================================================

\section{Change Detection}

\begin{figure}[H]
        \centering
        \includegraphics[width = 0.7\textwidth]{./Part1/figs/change_flow.png}
        \caption{Structure of the change detection based on hypothesis testing}
\end{figure}

Change detection or alarming is the task of finding abrupt or rapid changes in the estimate of the chosen parameter/state that is used as a guard condition to trigger another process like resetting the parameters of a filter or start the fault isolation and identification process, etc. This, the connecting step that takes the process from estimation to decision-making.

The problem of statistical change detection can be stated as the determination of change time $k_a$ or $t_a$ based on the predefined thresholds $h$ for the given time-series data $\theta_t$ with noise. This can be done both by using online (recursive) or offline (windowed) approaches. The change detection approaches can be divided into hypothesis testing and estimation/information approaches. Hypothesis testing algorithms have two distinct and independent parts to them, $1)$ A distance measure (or Residual) and $2)$ A stopping rule. In estimation approaches both of these parts are generally done in one step as the distance measure (residual) generated using this approach can be directly throsholeded without the averaging step of the stopping rule.

\subsection{Stopping Rule}

The stopping rule gives the alarm when $\hat \theta_t$ exceeds a certain threshold. The difference between the change detection and stopping rule is that the stopping rule captures the parameter exceeding the alarm which can be abrupt, linear or incipient (as in the case of aging detection or propeller drag is monotonically changing) but the change detection happens at a point of time when the decision needs to be made (catalyst has aged and is at EUL, propeller is broken). There is prior knowledge on the threshold for stopping rule. A fundamental nonlinear stopping rules is \cite{malladi1999generalized}:
\begin{itemize}
        \item Sequential Probability Ration Test (SPRT) algorithm
        \item Cumulative Sum (CUMSUM) algorithm (SPRT with no reset level)
\end{itemize}

\begin{algorithm}[H]
        \begin{align*}
                g_t &= g_{t-1} + y_t - \nu\\
                g_t &= 0 \: \text{and} \: \hat k = t \: \text{if} \: g_t < a < 0\\
                g_t &= 0 \: \text{and} \: t_a = t \: \text{and alarm if} \: g_t > h > 0\\
                \text{Design Parameter}&: \text{Drift } \nu, \text{ threshold } h, \text{ and reset level } a\\
                \text{Output}&: \text{Alarm time } t_a
        \end{align*}
        \caption{SPRT Algorithm (CUMSUM if a=0)}
\end{algorithm}

An application of CUMSUM algorithm for filtering test and truck data in presence of jumps in the process in presented in appendix. This test can be converted into a two-sided test by including another alarm condition with $-h$.


\subsection{Estimation/Information based approach to change detection (Distance Measure)}
Input to the stopping rule is a distance measure $s_t$ which can be the residual of the estimate or residual normalized to unit variance. The alternative is using the square of the residual (again, can be normalized to unit expectation of variance).

In estimation based approach the value of the distance measure can be directly used for generating alarm. Consider the linear parameter model:
\begin{align}
        y &= \phi^T \theta + \varepsilon
\end{align}
We want to test the following hypotheses:
\begin{itemize}
        \item[H0:] The parameter vectors are the same as nominal, $\theta = \theta_0$
        \item[H1:] The parameter vector $\theta$ is significantly different from the nominal $\theta_0$ and closer to $\theta_1$, the null hypothesis is rejected.
\end{itemize}
The distance measures that can be used to test the above hypotheses must not only include the detection for change in mean of $\hat \theta$ but also change in variance of $\hat \theta$. The following are a few distance measures from literature:

\begin{enumerate}
        \item \itbf{Prediction Error} \cite{segen1980detecting}:
        \begin{align}
                s^{PE} &= \frac{\norm{y - \phi^T \theta_0}}{\lambda_0} - N
        \end{align}
        \item \itbf{Generalized Likelihood Ratio} \cite{appel1983adaptive}:
        \begin{align}
                s^{GLR} &= N \log \frac{\lambda_0}{\lambda_1} + \frac{\norm{y - \phi^T \theta_0}^2}{\lambda_0}
                                                                - \frac{\norm{y - \phi^T \theta_0}^2}{\lambda_0}
        \end{align}
        \item \itbf{KL Divergence Test}\cite{kumamaru1996robust}:
        \begin{align}
                s^{DIV} &= N \lr{\frac{\lambda_0}{\lambda_1} - 1}
                                + \lr{1 + \frac{\lambda_0}{\lambda_1}} \frac{\norm{y - \phi^T \theta_0}^2}{\lambda_0}
                                - 2 \frac{\lr{y - \phi^T \theta_0}^T \lr{y - \phi^T \theta_1}}{\lambda_1}
        \end{align}
\end{enumerate}
$\lambda_0$, $\lambda_1$ are the nominal variance in the noise for the model with given parameters.

\section{Set Membership Identification}

Set membership identification is a class of techniques for estimating parameters of a linear or affine in parameter systems under a priori information about the error bounds that constraints the solution to certain set. These algorithms usually estimate discrete time affine-in-parameters models of the form:
\begin{align}
        y_t &= \theta^T x_t + \varepsilon_t
\end{align}
Where, $y_t$ is the output vector, $x_t$ is the sequence of measurement vectors (regressors) and $\varepsilon_t$ is an error or disturbance or an unknown input sequence which is magnitude bounded. Specifically, there exists a positive sequence $\lrf{\gamma_t}$ such that:
\begin{align}
        \abs{\varepsilon_t} < \gamma_t \qquad \forall \, t
\end{align}

Among the class of set membership techniques, Optimal Bounding Ellipsoid (OBE) algorithms are of particular interest for their superior adaptation, improved accuracy, efficient use of innovation data, improved computation efficiency, and increased robustness to measurement noise and to deviation from assumed noise model. OBE methods are algorithmically similar to recursive least squares but the philosophy, operation and behavior of the OBE methods is significantly different.

Geometrically, the error bounds at each time $t$ imply two hyperplanes in parameter space, between which the true value of the parameter vector must lie. These planes form an open 'hyper-strip' $H_t$ between the planes. These pairs of hyperplanes over time forms a sequence of polytopes $P_t$.
\begin{align}
        H^+_t &= \lrf{\theta | y_t = \theta^T x_t + \gamma_t} \qquad
        H^-_t  = \lrf{\theta | y_t = \theta^T x_t - \gamma_t}\\
        H_t   &= \lrf{\theta | \abs{y_t - \theta^T x_t} < \gamma_t}\\
        P_t   &= \bigcap_{t=1}^{\infty} H_t
\end{align}
The OBE algorithm produces a sequence of bounding hyper-ellipsoids, $\lrf{\mathcal{E}_t}_{t = m}^{\infty}$, which attempt, at each t, to tightly bound the exact polytopes in some sense. The interior of the hyper-ellipsoid at a time $t \geq m$ is the hyper-ellipsoidal membership set.
\begin{align}
        \mathcal{E}_t = \lrf{\theta | \lr{\theta - \theta_t}^T \kappa_t^{-1} C_t \lr{\theta - \theta_t} < 1}
\end{align}
$\theta_t$ is the centroid of the ellipsoid, often used as the point estimate of $\theta$ when required. The ellipsoids
are monotonically non-increasing in some measure of sizes and under specific conditions they converge to point
estimates. The parameters that describe the ellipsoid $\theta_t, \kappa_t, C_t = P_t^{-1}$ are recursively computed as follows:

\begin{algorithm}
        \begin{align*}
                \Gamma_t &= x_t^T P_{t-1} x_t \qquad (scalar)\\
                \varepsilon_t &= y_t - \theta_{t-1}^T x_t\\
                P_t &= \frac{1}{\alpha_t} \lrb{P_{t-1} - \frac{\beta_t P_{t-1} x_t x_t^T P_{t-1}}{\alpha_t + \beta_t \Gamma_t}}\\
                \theta_t &= \theta_{t-1} + \beta_t P_t x_t \varepsilon_t\\
                \kappa_t &= \alpha_t \kappa_{t-1} + \beta_t \gamma_t - \frac{\alpha_t \beta_t \varepsilon_t^2}{\alpha_t + \beta_t \Gamma_t}
        \end{align*}
        \caption{General framework of OBE algorithm's recursion \cite{deller1994unifying}}
\end{algorithm}

$\alpha_t$ and $\beta_t$ are positive weighing sequences chosen according to the particular OBE algorithm employed. The weights are chosen to explicitly or implicitly to minimize the 'size' of the set $\mathcal{E}_t$ at each iteration. When the optimal weights do not exist, the updating is stopped by setting $\beta_t  = 0$ and/or $\alpha_t = 1$.  $\alpha_t$ and $\beta_t$ can be tied together by making them a function of a single parameter $\lambda_i$. Commonly, there are three optimization criteria that can be employed:
\begin{enumerate}
        \item Determinant of the inverse of ellipsoid matrix which is proportional to the square of the volume of the ellipsoid.
        \begin{align}
                \mu_t^{vol} &= \det \lrf{\kappa_t C_t^{-1}}
        \end{align}

        \item Trace of the inverse ellipsoid matrix which is proportional to the sum of squares of the ellipsoid's semi-axes.
        \begin{align}
                \mu_t^{tr} &= trace \lrf{\kappa_t C_t^{-1}}
        \end{align}

        \item The parameter $\kappa_t$ (QOBE algorithm \cite{deller1993least}). This criterion in some cases results in directly ensuring the $\theta$ estimates lie in the hyperplanes.
\end{enumerate}
The general framework and the specific algorithms are presented in \cite{deller2002set}. The recursion part of the OBE algorithm is very similar to the recursive least-squares algorithm \cite{deller1989set} except for the choice of $\alpha_t$ and $\beta_t$ which are the result of minimizing the 'size' of the ellipsoid.



\subsection{Fault detection using SMI techniques}

The robustness of SMI to model structure error and noise makes it an ideal candidate for model-based fault diagnosis
using parameter estimates for reduced order models. \cite{combastel2016set} presents recent advances in application of
set membership techniques for fault diagnosis and fault tolerance in diverse systems. From change detection point of
view, the main techniques to detect parameter changes is using consistency tests that check if the intersection of
estimated ellipsoids with the nominal feasible parameter set is non-empty \cite{watkins1996fault},
\cite{ingimundarson2009robust}. This method can be further refined for isolation by including data-hypersector in which
the nominal parameter set lie but not the faulty parameter set \cite{reppa2016fault}.

\subsection{SMI in Adaptive Robust Control Structure}

The parameter bounds estimated using SMI techniques can improve the performance in Adaptive Robust Control designs as
the conservative estimates of the parameter bounds can be further reduced using ellipsoidal sets obtained recursively.
Such an approach of bound shrinking by using SMI techniques along with the usual parameter estimation using gradient
based adaptation low is presented in \cite{lu2009set} and input design for parameter estimation is presented in
\cite{lu2010experimental}. Further, \cite{kosut1992set} presents an adaptive control design using the combination of SMI
techniques and linear robust control design for linear systems. Proposed future work includes directly using centroid of
the ellipsoid for model compensation and the ellipsoidal bounds for projection in Adaptive Robust Control designs
avoiding the additional step of computing the gradient based parameter estimation.


\chapter{Contributions and Future Work}
The contributions of the present work include development, system identification and validation of diagnostic (and control) models for motor-propeller actuator system and diesel engine SCR system. These models are used for fault detection through parameter estimation through set membership techniques. A brief survey of methods of change detection and set membership based parameter estimation are also presented.

\subsection{Contributions in actuator fault detection and reconfiguration}
\begin{enumerate}
        \item Development and hardware implementation of angular velocity measurement system for motor propeller system using the noisy commutation signals from the motor.
        \item Development of the model of the actuator considering the high-frequency unknown nonlinear filters in the ESC of the actuator.
        \item Identification of the model using linearized model parameter estimation and validation of the nonlinear model across the full range operation (ACC'24 paper \cite{seshaIdent}).
        \item Closed-loop control of the actuator and preliminary work on fault detection. (MECC'24 (presented), SSRR'24 (presentedt))
\end{enumerate}

\subsection{Contributions in SCR-ASC aging diagnostics}
\begin{enumerate}
        \item Development, identification and validation of a linear diagnostic model using the CSTR and reduced order modelling assumptions. Demonstrating the viability of aging detection under limited operating range and absence of cross-sensitivity (MECC'24 (presented)).
        \item Development of a discrete nonlinear model using PFR (Plug Flow Reactor model) that explicitly considers the effects of interplay between sampling and residence time. The model also avoids the causality reversal that happens in the CSTR model where the rate constant is a function of tail-pipe concentrations instead of inlet concentrations.
        \item Development of an identifiable linear parameter model from the discrete nonlinear model for tailpipe $NO_x$ that is validated through parameter identification using only the available measurements on road conditions (No ammonia measurements).
\end{enumerate}

\section{Future Work and Publication Plan}
The future work includes development of set membership based parameter estimation schemes and integrating them with change detection methods for fault diagnostics. Then, validating the above method using the experimental data. In the case of hexrotor the future work also involves developing the reconfiguration system (Control for static hover under actuator failure) and demonstrating it using simulation and experiments.

\subsection{Set Membership based parameter estimation techniques}
\begin{enumerate}
        \item Integrating set membership methods with appropriate change detection methodology that can detect appropriate changes in parameter corresponding to specific actuator faults and SCR-ASC aging (Spring '25).
        \item Integrating set membership methods with adaptive robust control theory (Spring '25).
\end{enumerate}

\subsection{Actuator fault detection and reconfiguration}
\begin{enumerate}
        \item Demonstrating the performance gains in the actuator using RPM feedback and SMI based adaptive robust control. (Journal paper 1 in IJRR. Draft by January'25)
        \item Developing and demonstrating the actuator fault detection that detects propeller damage, breaking and touching surfaces using set membership techniques. (Journal Paper 2, Spring '25)
        \item Developing the reconfiguration of the tilted hexrotor into static hover under actuator faults while compensating for the unbalanced angular momentum using RPM feedback. Demonstrating the control design using simulation and with available resources an experimental setup (Journal Paper 3. Summer '25 and beyond)
\end{enumerate}

\subsection{SCR-ASC aging diagnostics}
\begin{enumerate}
        \item Demonstrating the validity of the parametric model on truck data. Publishing (with approval from Cummins) the model development and validation results. (Journal paper 4. JDSMC. February '25)
        \item Developing a fault detection methodology, based on the change in specific parameter sets, for the nonlinear model and validating it with the experimental data. (Journal paper 5. JDSMC. Spring&Summer '25).
        \item Implementation considerations of the algorithm using available hardware.
\end{enumerate}

